\documentclass[11pt]{article}
\usepackage{mathunal-base}
\mathunalfuente{serif}

\renewcommand{\examtitulo}{Métodos Numéricos 3002192 --- Segundo Parcial}
\renewcommand{\examfecha}{Octubre 7 de 2006}

\begin{document}

\examheader

\vspace{6pt}
\noindent Nombre \dotfill\ Carné \dotfill\ Grupo \dotfill

\vspace{10pt}
\noindent\textbf{1. (15\%)} La tabla de diferencias divididas para el polinomio interpolante de una función $f(x)$ definida en $[0,0.7]$ está dada por
\[
\begin{array}{c|c|c|c}
x_0=0 & f[x_0] & & \\ \hline
x_1=0.4 & f[x_1] & f[x_0,x_1] & \\ \hline
x_2=0.7 & f[x_2]=6 & f[x_1,x_2]=10 & f[x_0,x_1,x_2]=\dfrac{50}{7}
\end{array}
\]
Encuentre $f[x_0]$, $f[x_1]$ y $f[x_0,x_1]$.

\vspace{5mm}
\noindent\textbf{2. (15\%)} Una cercha cúbica sujeta $S$ para una función $f(x)$ definida en $[1,3]$ está dada por
\[
S(x)=\begin{cases}
S_0(x)=3(x-1)+2(x-1)^2-(x-1)^3, & 1\leq x\leq 2\\
S_1(x)=a+b(x-2)+c(x-2)^2+d(x-2)^3, & 2\leq x\leq 3
\end{cases}
\]
Si $f'(1)=f'(3)$, encuentre $a$, $b$, $c$ y $d$.

\vspace{5mm}
\noindent\textbf{3.} La recta de regresión, que es la que mejor se aproxima en mínimos cuadrados a una colección de datos $\{(x_j,y_j)\}_{j=1}^{N}$, es $y=Ax+B$. Sus coeficientes $A$ y $B$ son la solución del sistema de ecuaciones normales
\[
\left(\sum_{k=1}^N x_k^2\right)A+\left(\sum_{k=1}^N x_k\right)B=\sum_{k=1}^N x_ky_k
\]
\[
\left(\sum_{k=1}^N x_k\right)A+NB=\sum_{k=1}^N y_k
\]
\begin{enumerate}
\item[a.] (15\%) Sean $\bar x=\dfrac1N\displaystyle\sum_{k=1}^N x_k$ y $\bar y=\dfrac1N\displaystyle\sum_{k=1}^N y_k$ las medias aritméticas para los puntos $\{(x_j,y_j)\}_{j=1}^N$. Pruebe que el punto $(\bar x,\bar y)$ está en la recta de regresión determinada por el conjunto de puntos dado.
\item[b.] (15\%) Encuentre $A$ y $B$, los coeficientes de la recta de regresión, sabiendo que los datos están dados por la siguiente tabla
\[
\begin{array}{c|c|c|c}
x_j & 2 & 3 & 4\\ \hline
y_j & \frac15 & \frac1{10} & \frac1{17}
\end{array}
\]
\end{enumerate}

\vspace{5mm}
\noindent\textbf{4. (20\%)} Determine las constantes $w_1$, $x_0$ y $x_1$ de manera que la regla de cuadratura
\[
\int_0^1 f(x)\,dx=\frac12 f(x_0)+w_1f(x_1)
\]
sea exacta para todos los polinomios de grado $\leq 3$.

\vspace{5mm}
\noindent\textbf{5. (20\%)} Se utiliza el método de Romberg para aproximar $\displaystyle\int_2^3 f(x)\,dx$. La tabla es:
\[
\begin{array}{l}
R(0,0)\\
R(1,0)\quad R(1,1)\\
R(2,0)\quad R(2,1)\quad R(2,2)
\end{array}
\]
Se sabe que $f(2)=0.51342$, $f(3)=0.36788$, $R(2,0)=0.43687$ y $R(2,2)=0.43662$. Encuentre $f(2.5)$.

\vfill
\rule{\linewidth}{0.4pt}\\[1pt]
{\scriptsize\textit{Transcripción digital --- MathUNAL}\hfill\textcolor{unalblue}{\textbf{1}}}

\end{document}
